% !TEX root = ../../proposal.tex

\chapter{Conclusion}
\label{chap:conclusion}

This proposal argues that the main barrier to reliable compositional generation is the mismatch
between logical composition and the semantically correct joint concept. To address this problem,
the project proceeds in three stages: first, it characterises the composability gap
through a structured benchmark and trajectory-based analysis; second, it tests whether more
factorised latent representations can reduce this gap in a controlled setting; and third, it
evaluates whether the same principles can support clinically plausible synthesis of rare
co-morbid chest X-rays.

Taken together, the proposed work aims to contribute both a clearer theoretical account of why
compositional diffusion models fail and a practical pathway for improving out-of-distribution
generation in low-resource domains. If successful, the thesis will show that reliable composition
depends not only on how conditions are combined at inference time, but also on how cleanly the
underlying concepts are represented, with direct implications for generative modeling and medical
image synthesis.
